\chapter{在磁场中的原子}
\problem{6.1}{已知钒原子的基态是 ${}^4\mathrm F_{3/2}$。
\begin{enumerate}
\item 问钒原子束在不均匀横向磁场中将分裂为几束？
\item 求基态钒原子的有效磁矩 $\mu_J$。
\end{enumerate}}
\answerarea
\problem{6.2}{已知 $\mathrm{He}$ 原子 ${}^1\mathrm P_1\to{}^1\mathrm S_0$ 跃迁的光谱线在磁场中分裂为三条光谱线，其间距 $\Delta\widetilde\nu=0.467\,\mathrm{cm^{-1}}$，试计算所用磁场的磁感应强度。}
\answerarea\clearpage
\problem{6.3}{$\mathrm{Li}$ 漫线系的一条（$3\,{}^2\mathrm D_{3/2}\to2\,{}^2\mathrm P_{1/2}$）在弱磁场中将分裂成多少条光谱线？试作出相应的能级跃迁图。}
\answerarea
\problem{6.4}{在平行于磁场方向观察到某光谱线的正常塞曼效应分裂的两谱线间波长差为 $\Delta\lambda=0.40\,\angstrom$。所用的磁场的 $B$ 是 $2.5\,\mathrm T$。试计算该谱线原来的波长。}
\answerarea\clearpage
\problem{6.5}{氦原子光谱中波长为 $6678.1\,\angstrom$（$1\mathrm s3\mathrm d\,{}^1\mathrm D_2\to1\mathrm s2\mathrm p\,{}^1\mathrm P_1$）及 $7065.1\,\angstrom$（$1\mathrm s3\mathrm s\,{}^3\mathrm S_1\to1\mathrm s2\mathrm p\,{}^3\mathrm P_0$）的两条谱线，在磁场中发生塞曼效应时各分裂成几条？分别作出能级跃迁图。问哪一个是正常塞曼效应？哪一个不是？为什么？}
\answerarea
\problem{6.6}{$\mathrm{Na}$ 原子从 $3\,{}^2\mathrm P_{1/2}\to3\,{}^2\mathrm S_{1/2}$ 跃迁的光谱线波长为 $5896\,\angstrom$，在 $B=2.5\,\mathrm T$ 的磁场中发生塞曼分裂。问从垂直于磁场方向观察，其分裂为多少条光谱线？其中波长最长和最短的两条光谱线的波长各为多少？}
\answerarea\clearpage
\problem{6.7}{$\mathrm{Na}$ 原子 $3\mathrm P\to3\mathrm S$ 跃迁的精细结构为两条，波长分别为 $5895.93\,\angstrom$ 和 $5889.96\,\angstrom$。试求出原能级 ${}^2\mathrm P_{3/2}$ 在磁场中分裂后的最低能级与 ${}^2\mathrm P_{1/2}$ 分裂后的最高能级相合时所需要的磁感应强度 $B$。}
\answerarea
\problem{6.8}{已知铁（${}^5\mathrm D$）的原子束在横向不均匀磁场中分裂为 9 束。问铁原子的 $J$ 值多大？其有效磁矩多大？如果已知上述铁原子的速度 $v=10^3\,\mathrm{m/s}$，铁原子量为 55.85，磁极范围 $L_1=0.03\,\mathrm m$（习题图 6.1），磁铁到屏的距离 $L_2=0.10\,\mathrm m$，磁场中横向的磁感应强度的不均匀梯度 $\mathrm dB/\mathrm dy=10^3\,\mathrm{T/m}$，试求屏上偏离最远的两束之间的距离 $d$。
\fig{6_1}{6.1}}
\answerarea\clearpage
\problem{6.9}{铊原子气体在 ${}^2\mathrm P_{1/2}$ 状态。当磁共振时 $B=0.2\,\mathrm T$，观察到顺磁共振现象。问微波发生器的频率多大？}
\answerarea
\problem{6.10}{钾原子在 $B=0.3\,\mathrm T$ 的磁场中，当交变电磁场的频率为 $8.4\times10^9\,\mathrm{Hz}$ 时观察到顺磁共振。试计算朗德因子 $g$，并指出原子处在何种状态？}
\answerarea\clearpage
\problem{6.11}{氩原子（$Z=18$）的基态为 ${}^1\mathrm S_0$；钾原子（$Z=19$）的基态为 ${}^2\mathrm S_{1/2}$；钙原子（$Z=20$）的基态为 ${}^1\mathrm S_0$；钪原子（$Z=21$）的基态为 ${}^2\mathrm D_{3/2}$。问这些原子中哪些是抗磁性的？哪些是顺磁性的？为什么？}
\answerarea\clearpage
